\section{Limitations and Threats to Validity}
\label{sec:limitations}

While the experimental results presented in this study offer promising evidence for the feasibility of LLM-driven, constraint-based MVC code generation, several limitations and threats to validity must be acknowledged. Following established practice in empirical software engineering research, these threats are organized across four validity dimensions.

\paragraph{Threats to External Validity:}
External validity concerns the extent to which the findings generalize beyond the specific conditions of this study. The evaluation corpus comprises a purposive, criterion-based sample of 20 SRS documents drawn from high-performing senior-level software engineering capstone projects (scoring 90/100 or above), substantially expanding the five-document corpus of our preliminary study \citep{dim2026mvc}. This sample is substantial for a qualitative, expert-validated evaluation of this kind, though, as with most case-study-based empirical software engineering research, a larger corpus would further strengthen confidence in the stability of the reported performance figures. While these documents represent well-structured, high-fidelity specifications, they originate from an academic setting and predominantly cover web-based management and reservation systems. The system's performance on documents from domains with substantially different structural conventions, such as embedded systems, real-time control applications, or safety-critical software, has not been assessed; consequently, the reported results should not be generalized to industrial, safety-critical, embedded, or poorly specified projects. Furthermore, the current implementation is designed exclusively for the MVC architectural pattern and does not directly generalize to other paradigms, such as microservices, event-driven architectures, or n-tier systems, without substantial re-engineering. A further threat concerns the fault-injection protocol itself (Section~\ref{sec:experiment}): the 406 violations used to evaluate detection and repair performance were deliberately and synthetically introduced across four predefined rule categories, and while this design enables precise ground-truth measurement, it may not fully capture the diversity, subtlety, or co-occurrence patterns of violations that arise organically during real-world, iterative development. Consequently, the reported VDR, ACP, RSDA, and RSR\textsubscript{e2e} figures should be interpreted as performance on a controlled benchmark rather than a guarantee of equivalent effectiveness against naturally occurring architectural drift.

\paragraph{Threats to Construct Validity:}
Construct validity concerns whether the chosen metrics accurately represent the theoretical concepts they are intended to measure. The current prototype caps the output at eight Model entities, six Controllers, and six primary Views to bound prompt and artifact size; these caps reflect an engineering trade-off rather than a claim about the ideal number of components for any given system, and they may truncate valid components in larger specifications, potentially reducing extraction completeness. BERTScore, while widely adopted for semantic similarity evaluation, serves as a proxy for semantic fidelity rather than a direct measure of requirement completeness. High BERTScore values indicate strong lexical and contextual similarity but do not imply that all functional requirements have been correctly captured. Similarly, the ACES metric aggregates expert scores across a structured rubric to represent architectural viability; however, the degree to which expert agreement on a 0--5 Likert scale reflects real-world industrial standards remains an assumption of the evaluation design.

\paragraph{Threats to Internal Validity:}
Internal validity concerns whether the observed outcomes can be attributed to the system's design rather than confounding factors. The ACES metric relies on evaluations conducted by a panel of five expert software engineers (N=5). While this panel size is consistent with established heuristic evaluation practice~\citep{monasor2014walk}, individual differences in experience, domain background, and interpretation of the rubric criteria may introduce subjectivity into the scores. To mitigate this threat, an Inter-Rater Reliability analysis was conducted (ICC(A,5)\,$=$\,0.942, 95\% CI [0.93, 0.95]; see Section~\ref{sec:discussion_rq4} for the full methodology), indicating excellent agreement among the five expert raters and supporting the reliability of the reported ACES scores; however, the influence of rubric design on scoring behavior cannot be entirely ruled out. This study reports a within-system, descriptive evaluation of one implementation of the \textit{MVC Test Orchestrator}. Because no matched baseline or component ablation was evaluated on the same SRS corpus, the observed results cannot be attributed to the multi-agent decomposition, RAG grounding, or rule-based stages individually. The findings should therefore be interpreted as evidence of feasibility under the tested configuration rather than as evidence of superiority or causal component effectiveness.

\paragraph{Threats to Conclusion Validity:}
Conclusion validity concerns the reliability of the relationships drawn from the data. The three-iteration consistency analysis (RQ5) provides a practical approximation of run-to-run variance; however, three iterations represent a limited sample for characterizing the full stochastic behavior of the underlying model. Additionally, the evaluation was conducted using a fixed version of \textit{Gemini 2.5 Flash} with a specific set of hyperparameters. Different model versions, API updates, or alternative temperature settings may produce different consistency and remediation outcomes, limiting the stability of the reported conclusions across model releases. Furthermore, the HITL checkpoints were not exercised during the reported experiments; their effect on output quality, time, or developer acceptance remains unquantified.
